% Generated draft: requires mathematical and source review.
\documentclass[11pt]{article}
\usepackage[margin=25mm]{geometry}
\usepackage{fontspec}
\setmainfont{Latin Modern Roman}
\usepackage{amsmath,amssymb,mathtools,unicode-math}
\setmathfont{Latin Modern Math}
\usepackage{enumitem,needspace,xurl,hyperref}
\setlength{\parindent}{0pt}
\setlength{\parskip}{5pt}
\setlength{\emergencystretch}{3em}
\newcommand{\N}{\mathbb{N}}
\newcommand{\Z}{\mathbb{Z}}
\newcommand{\Q}{\mathbb{Q}}
\newcommand{\R}{\mathbb{R}}
\newcommand{\C}{\mathbb{C}}
\newcommand{\F}{\mathbb{F}}
\newcommand{\A}{\mathbb{A}}
\newcommand{\PP}{\mathbb{P}}
\newcommand{\E}{\mathbb{E}}
\newcommand{\eps}{\varepsilon}
\newcommand{\dd}{\,\mathrm{d}}
\begin{document}
\section*{UnsolvedMath ID: 800709}
\textbf{Source problem code:} catalogue entry
\subsection{Definitions and mathematical statement}
Let \(F\) be a field and let \(A\) be a finite-dimensional central simple \(F\)-algebra, meaning a nonzero associative unital \(F\)-algebra with no nontrivial two-sided ideals and with center \(F\). An involution of the first kind on \(A\) is an \(F\)-linear map \(\sigma:A\to A\) satisfying
\[
\sigma(ab)=\sigma(b)\sigma(a),\qquad \sigma^2=\operatorname{id}_A.
\]
After extending scalars to an algebraic closure of \(F\), such an involution is adjoint to a nonsingular bilinear form: the adjoint relation is \(b(Tu,v)=b(u,\sigma(T)v)\). The involution is called \emph{symplectic} when this form is alternating, and \emph{orthogonal} when it is symmetric and nonalternating. Here alternating means \(b(v,v)=0\) for every vector \(v\).

The involution \(\sigma\) is \emph{anisotropic} if
\[
\forall x\in A,\qquad \sigma(x)x=0\ \Longrightarrow\ x=0.
\]
For a field extension \(K/F\), set \(A_K=A\otimes_F K\) and \(\sigma_K=\sigma\otimes\operatorname{id}_K\).

\textbf{Question (supplied Problem 19).} For every field \(F\), every central simple \(F\)-algebra \(A\), every anisotropic orthogonal or symplectic involution \(\sigma\) on \(A\), and every finite field extension \(K/F\) of odd degree, must \(\sigma_K\) be anisotropic? Equivalently, does
\[
\left(\forall x\in A,\ \sigma(x)x=0\Longrightarrow x=0\right)
\quad\Longrightarrow\quad
\left(\forall y\in A_K,\ \sigma_K(y)y=0\Longrightarrow y=0\right)
\]
hold whenever \([K:F]\) is odd?

\emph{Scope qualification.} The supplied statement gives no characteristic restriction and no separability hypothesis on \(K/F\); none is added here. Thus this rendering includes characteristic \(2\) and finite inseparable extensions of odd degree where applicable. Whether the source list imposes additional ambient conventions is unresolved from the supplied record.
\subsection{Short English statement}
Does every anisotropic orthogonal or symplectic involution on a central simple algebra remain anisotropic over every finite field extension of odd degree? The supplied statement imposes no characteristic or separability restriction; any additional conventions in the source list require review.
\subsection{Sources}
\begin{itemize}
\item Local UnsolvedMath registry: permanent ID \(800709\), exact source problem code \texttt{catalogue entry}, title \emph{Open Problems on Central Simple Algebras --- Problem 19}. The registry attributes the problem to Auel, Brussel, Garibaldi, and Vishne, with proposed year 2010.
\item Auel, Brussel, Garibaldi, and Vishne, \emph{Open Problems on Central Simple Algebras} (2010), Problem 19. Supplied source URL: \url{https://arxiv.org/abs/1006.3304}. The supplied provenance records access on 2026-10-16 and extraction method \texttt{arxiv\_latex}; these are source-record assertions, not a claim of independent access or verification here.
\item The local registry labels the problem \emph{open}. This is supplied curation, not an independently verified assessment of the literature or current mathematical status.
\end{itemize}
\end{document}
